\subsection{Ryll-Nardzewski theorem}

In this section, E denotes a normed space over the field R and T a Hausdorff locally convex topology on E for which the norm of E is lower semi-continuous. These hypotheses are in particular satisfied in the following cases :

\begin{enumerate}
\def\labelenumi{\alph{enumi})}
\item
  \(\mathcal{T}\) is the topology induced by the norm of the normed space E.
\item
  \(\mathcal{T}\) is the weakened topology \(\sigma (\mathrm{E},\mathrm{E}^{\prime})\) of the normed space E.
\item
  E is the dual of a normed space F and \(\mathcal{T} = \sigma (\mathrm{F}',\mathrm{F})\)
\end{enumerate}

\(d\) ) There exist two normed spaces \(\mathrm{F}_1\) and \(\mathrm{F}_2\) such that \(\mathrm{E} = \mathcal{L}(\mathrm{F}_1;\mathrm{F}_2)\) and \(\mathcal{T}\) is the topology of simple convergence.

Unless otherwise expressly stated, the topological notions refer to the topology \(\mathcal{T}\) .

Let K be a convex subset of E. Suppose that K is compact (for the topology T), and that K satisfies the first axiom of countability for the distance defined by the norm of E.

Lemma 2. --- Suppose K contains at least two points. For every \(\varepsilon > 0\) , there exists a partition of K into two non-empty subsets \(\mathbf{K}_1\) and \(\mathbf{K}_2\) , having the following properties:

\begin{enumerate}
\def\labelenumi{\alph{enumi})}
\item
  \(\mathrm{K}_1\) is convex and compact;
\item
  we have \(\| x_1 - x_2\| < \varepsilon\) for every \(x_{1}\) and \(x_{2}\) in \(\mathbf{K}_2\) .
\end{enumerate}

Let L be the closure of the set of all extremal points of K. By the Krein-Milman th. (II, p.~55, th. 1), K is the closed convex envelope of L. Since K contains at least two points, so does L. For every \(x \in L\) , let \(A_{x}\) be the set of all \(y \in L\) such that \(\|x - y\| \leqslant \varepsilon/4\) . By the hypothesis, on K, there exists a countable subset D of L such that \(L = \bigcup_{x \in D} A_{x}\) . Since the norm is lower semi-continuous, each of the sets

\(A_{x}\) is closed. Applying Baire's th. (GT, IX, § 5, No.~3, th. 1) to the compact space L, we see that there exists a point a in D and an open subset U in E such that \(L \cap U\) is non-empty and is contained in \(A_{a}\) . Since L contains at least two points, and since E is Hausdorff, we can choose U in such a way that \(L \not\subset U\) .

Let M be the closed convex envelope of \(L \cap \complement U\) . For every real number t such that 0 \textless{} t \textless{} 1, let \(M_{t}\) be the set of all vectors of the form \(tx_{1} + (1 - t)x_{2}\) with \(x_{1} \in M\) and \(x_{2} \in K\) ; this is a non-empty, compact convex subset of K. We shall prove that \(M_{t} \neq K\) by reductio ad absurdum. Suppose that \(M_{t} = K\) ; then every extremal point x in K belongs to \(M_{t}\) , hence can be written in the form \(x = tx_{1} + (1 - t)x_{2}\) with \(x_{1} \in M\) and \(x_{2} \in K\) . This implies that \(x = x_{1} = x_{2}\) , and so \(x \in M\) . By Krein-Milman th. (II, p.~55, th. 1), we have K = M, and K is the closed convex envelope of \(L \cap \complement U\) . By II, p.~56, corollary, this implies that \(L \subset L \cap \complement U\) , which contradicts the relation \(L \cap U \neq \varnothing\) .

Put \(d = \sup_{x \in \mathbf{K}, y \in \mathbf{K}} \| x - y \|\) and choose a real number \(t\) such that \(0 < t < 1\) and \(t < \varepsilon / 4d\) . Put \(\mathbf{K}_1 = \mathbf{M}_t\) and \(\mathbf{K}_2 = \mathbf{K} - \mathbf{M}_t\) . By the preceding argument, the sets \(\mathbf{K}_1\) and \(\mathbf{K}_2\) are non-empty, and \(\mathbf{K}_1\) is convex and compact. Let \(\mathbf{M}'\) be the closed convex envelope of \(\mathbf{L} \cap \mathbf{U}\) . Since \(\mathbf{K}\) is the closed convex envelope of the set \(\mathbf{L} = (\mathbf{L} \cap \complement \mathbf{U}) \cup (\mathbf{L} \cap \mathbf{U})\) , it is also the closed convex envelope of \(\mathbf{M} \cup \mathbf{M}'\) . Let \(x_1\) and \(x_2\) be two points in \(\mathbf{K}_2\) ; for \(i = 1, 2\) , there exist \(y_i \in \mathbf{M}\) , \(z_i \in \mathbf{M}'\) and a real number \(\alpha_i\) such that \(0 \leqslant \alpha_i \leqslant 1\) and \(x_i = \alpha_i y_i + (1 - \alpha_i) z_i\) . If \(\alpha_i \geqslant t\) , then \(x_i = ty_i + (1 - t) \left\{\frac{\alpha_i - t}{1 - t} y_i + \frac{1 - \alpha_i}{1 - t} z_i\right\}\) ; this contradicts the assumption that \(x_i \notin \mathbf{M}_t\) . Hence \(\alpha_i < t\) , for \(i = 1, 2\) , and so

\[
\left\| x _ {i} - z _ {i} \right\| = \left\| \alpha_ {i} \left(y _ {i} - z _ {i}\right) \right\| = \alpha_ {i} \left\| y _ {i} - z _ {i} \right\| \leqslant \alpha_ {i} d <   d t <   \varepsilon / 4.
\]

For every point \(z\) in \(\mathbf{M}'\) , we have \(\| z - a \| \leqslant \varepsilon / 4\) , since \(\mathrm{L} \cap \mathrm{U} \subset \mathrm{A}_a\) , and so, in particular \(\| z_i - a \| \leqslant \varepsilon / 4\) . Thus

\[
\left\| x _ {1} - x _ {2} \right\| \leqslant \sum_ {i = 1} ^ {2} \left(\left\| x _ {i} - z _ {i} \right\| + \left\| z _ {i} - a \right\|\right) <   \varepsilon .
\]

This completes the proof.

Lemma 3. --- Let G be a group of continuous (for T) affine transformations on K. Suppose that K is non-empty and that \(\|gx - gy\| = \|x - y\|\) for all x, y in K and all g in G. Then there exists a point in K which is invariant under G.

Let \(\mathfrak{J}\) be the family of non-empty subsets of \(K\) which are closed convex and stable for \(G\) . If \((L_{\alpha})_{\alpha \in I}\) is a family of elements of \(\mathfrak{J}\) which is totally ordered by inclusion, then the set \(L = \bigcap_{\alpha \in I} L_{\alpha}\) belongs to \(\mathfrak{J}\) . Consequently (S, III, § 3, No.~4, th. 2), there exists

an element \(\mathbf{L}\) in \(\mathfrak{J}\) which is minimal for the relation of inclusion. We shall prove that \(\mathbf{L}\) reduces to a point.

We argue by reductio ad absurdum, assuming that L contains at least two distinct points \(x_{1}\) and \(x_{2}\) , put \(x = (x_{1} + x_{2})/2\) and \(\varepsilon = \|x_{1} - x_{2}\|/2\) . The convex and compact set L satisfies the first axiom of countability for the distance defined by the norm (GT, IX, § 2, No.~8). Hence we can apply lemma 2 and find a compact and convex subset \(L_{1}\) of L, distinct from \(\varnothing\) and from L, having the following property :

\begin{enumerate}
\def\labelenumi{(\Alph{enumi})}
\tightlist
\item
  For every \(y_{1}\) and \(y_{2}\) in \(L - L_{1}\) , we have \(\|y_{1} - y_{2}\| < \varepsilon\) .
\end{enumerate}

We shall prove by reductio ad absurdum that \(gx \in L_{1}\) for all \(g \in G\) . Let \(g \in G\) be such that \(gx \in L - L_{1}\) then we have

\[
\left\| g x _ {i} - g x \right\| = \left\| x _ {i} - x \right\| = \left\| x _ {1} - x _ {2} \right\| / 2 = \varepsilon ,
\]

for \(i = 1, 2\) . By property (A), we have \(gx_{i} \in \mathbf{L}_{1}\) . Since \(\mathbf{L}_{1}\) is convex, we conclude that \(gx = (gx_{1} + gx_{2}) / 2\) belongs to \(\mathbf{L}_{1}\) , which contradicts the assumption.

Let \(L'\) be the closed convex envelope of the orbit Gx of x. The set \(L'\) belongs to J. By the preceding argument, we have \(L' \subset L_1\) , hence \(L' \subset L\) , \(L' \neq L\) . This contradicts the minimal character of L and the proof is complete.

THEOREM 2 (Ryll-Nardzewski). --- Let E be a normed space and K a non-empty convex subset of E, which is compact for the weakened topology \(\sigma(E, E')\) . Let G be a group of isometric affine transformations of K. Then there exists a point in K which is invariant under G.

For every \(g \in \mathbf{G}\) , let \(\mathbf{K}_g\) denote the set of all points \(x\) in \(\mathbf{K}\) such that \(gx = x\) ; let \(\mathbf{K}\) be assigned the weakened topology; each set \(\mathbf{K}_g\) is convex and closed in the compact space \(\mathbf{K}\) . We shall prove that the intersection \(\bigcap_{g \in \mathbf{G}} \mathbf{K}_g\) is non-empty; for this, it is

enough to prove that the set \(K_{g_{1}} \cap \ldots \cap K_{g_{n}}\) is non-empty for every \(g_{1}, \ldots, g_{n}\) in G. Fix \(g_{1}, \ldots, g_{n}\) and let H be the subgroup of G generated by \(\{g_{1}, \ldots, g_{n}\}\) . Choose a point a in K and let L denote the closed convex envelope of the orbit Ha of a. Let D be the countable set of elements of the form \(\lambda_{1}h_{1}a + \cdots + \lambda_{m}h_{m}a\) , where \(\lambda_{1}, \ldots, \lambda_{m}\) are positive rational numbers such that \(\lambda_{1} + \cdots + \lambda_{m} = 1\) , and \(h_{1}, \ldots, h_{m}\) are elements in H. The closure \(\overline{D}\) of D for the strong topology, is convex, hence it is closed for \(\sigma(\mathrm{E}, \mathrm{E}')\) (IV, p.~4, prop. 2); therefore \(\overline{D} = L\) and this proves that L is a metric space satisfying the first axiom of countability for the distance \((x, y) \mapsto \|x - y\|\) . We can now apply lemma 2. There exists a point b in L which is invariant under H, hence \(b \in K_{g_{1}} \cap \ldots \cap K_{g_{n}}\) .

COROLLARY. --- Let E be a reflexive Banach space, G a group of automorphisms of the normed space E, and K a subset of E. Suppose that K is non-empty, convex, closed, bounded and stable under G. Then there exists a point in K which is invariant under G.

Since E is reflexive, K is compact for \(\sigma(\mathrm{E}, \mathrm{E}')\) (IV, p.~15, th. 1). Moreover, every element of G belongs to \(\mathcal{L}(\mathrm{E})\) .

